\begin{proofbox}
	\textbf{\textcolor{CProof}{思路提示}}

	\textbf{由于不等式两边均为非负，可以将两边平方后比较，转化为判断一个代数式的非负性。}

	\medskip
	\textbf{\textcolor{CProof}{正式推导}}

	\begin{description}[style=nextline, leftmargin=2.8em, labelsep=0.5em, itemsep=0.55em]
		\item[\textbf{步骤一（平方转化）：}]
		      因为 $\sqrt{u}$、$\sqrt{v}$、$\sqrt{2(u+v)}$ 均非负，故原不等式等价于
		      \[
			      (\sqrt{u}+\sqrt{v})^2 \le 2(u+v).
		      \]
		      \par\textbf{理由：} 对非负实数来说，$a\le b$ 与 $a^2 \le b^2$ 等价（平方函数在 $[0,+\infty)$ 单调递增）。

		\item[\textbf{步骤二（作差变形）：}]
		      计算右边与左边的差：
		      \[
			      \begin{aligned}
				      2(u+v) - (\sqrt{u}+\sqrt{v})^2
				       & = 2u + 2v - (u + v + 2\sqrt{uv}) \\
				       & = u + v - 2\sqrt{uv}             \\
				       & = (\sqrt{u} - \sqrt{v})^2.
			      \end{aligned}
		      \]
		      \par\textbf{理由：} 利用完全平方公式 $(\sqrt{u}+\sqrt{v})^2 = u + v + 2\sqrt{uv}$ 展开后合并同类项。

		\item[\textbf{步骤三（非负判断）：}]
		      对于任意实数，平方数都非负，故
		      \[
			      (\sqrt{u} - \sqrt{v})^2 \ge 0.
		      \]
		      \par\textbf{理由：} 实数的平方总是大于或等于零。当且仅当 $\sqrt{u}=\sqrt{v}$，即 $u=v$ 时等于零。

		\item[\textbf{步骤四（回代结论）：}]
		      由步骤一至三得到 $2(u+v) - (\sqrt{u}+\sqrt{v})^2 \ge 0$，即 $(\sqrt{u}+\sqrt{v})^2 \le 2(u+v)$。两边开方得原不等式 $\sqrt{u}+\sqrt{v} \le \sqrt{2(u+v)}$，且等号当且仅当 $u=v$ 时成立。
		      \par\textbf{理由：} 平方根函数在 $[0,+\infty)$ 上是单调递增的，从而平方不等式在开方后不等号方向不变。
	\end{description}

	\medskip
	\textbf{\textcolor{CProof}{结论回扣}}

	\textbf{至此，我们严格证明了对于任意非负实数 $u,v$，不等式 $\sqrt{u}+\sqrt{v} \le \sqrt{2(u+v)}$ 恒成立，且等号成立当且仅当 $u=v$。当 $u+v$ 为定值时，这一定量关系可直接用于含根式和的最大值计算。}
\end{proofbox}